\scanpage{227}
\mspage{116}
\[
 \Pi_Z^{\wedge}=\text{ens des }Z\text{-jauges }\deleted{\text{bijectives}}\text{ de }\Gamma
 \tag*{(25)}
\]
alors $Z$ opère sur $\Pi_Z^\wedge$ par
\[
 T_\lambda\varphi\overset{\mathrm{déf}}=\lambda_{S^\vee}\circ\varphi=\varphi\circ\lambda_S^{-1}
 \tag*{(25)}
\]
-- l'opération correspondante sur l'ens des graphes
associés aux $\varphi$ étant donnée par
\[
 T_\lambda A_\varphi=\{\{s,\lambda_S(t)\}\in\mathcal P(S)\mid s,t\ \text{t.q. }\{s,t\}\in A_\varphi\}.
 \tag*{(27)}
\]
Si $\varphi$ est surjective, alors pour $\varphi$ fixé, la connaissance
de $T_\lambda\varphi=\lambda_{S^\vee}\circ\varphi$ implique celle de $\lambda_{S^\vee}$, si $\varphi$ est injective
elle implique celles de $\lambda_S$, donc si $\varphi$ est bijective,
la connaissance de $T_\lambda\varphi$ implique celle de $\lambda$, i.e.
\deleted{\unsure{l'opérat|l'opération}}

(27) sur l'ens \inserted{$\Pi_Z$} des $Z$-jauges \uline{bijectives}, $Z$ opère
\uline{librement} -- donc sur \deleted{un graphe sé} \inserted{l'ens des}
structures de graphes \uline{séparés} sur un ens $S$,
admettant un sous-groupe comm. $Z_S\subset\unsure{\mathfrak S|\mathfrak G}_S$
comme groupe d'homothétie, $Z$ opère
librement.

Notons
\[
 G=\mathrm{Cent}_A(Z)\deleted{\unsure{s}},\text{ d'où }\quad Z\subset G\subset A
 \tag*{(28)}
\]
On trouve
\[
\left\{\begin{aligned}
 &\text{Soit }\varphi_0\in\Pi_Z\text{ une }Z\text{-jauge stricte }\inserted{\text{(i.e. bij)}}\text{ de }\Gamma,
 \ \deleted{\text{alors}}\inserted{\text{soit}}\\
 &\deleted{\unsure{G_{\varphi_0}|\mathcal G_{\varphi_0}}}\quad G_{\varphi_0}=\{g\in G\mid g_{S^\vee}\varphi=\varphi_{\unsure{0|\vphantom{0}}}g_S^{-1}
 \quad\text{i.e. }\varphi_0\text{ est une }g\text{-jauge}\}\\
 &(\text{\uline{NB} }G_{\varphi_0}\text{ n'est pas néc. un sous-groupe de }G)\\
 &\text{alors}\qquad\deleted{\text{alors}}\quad\deleted{\unsure{(30)}}\\
 &g\longmapsto\deleted{\unsure{g}}T_g\varphi\overset{\mathrm{déf}}=g_{S^\vee}\varphi_0=\varphi_0g_S^{-1}\\
 &\text{est une bijection}\\
 &G_{\varphi_0}\xrightarrow{\sim}\Pi_Z
\end{aligned}\right.\tag*{(29)}
\]
Mais on fera attention que par (23), $G_{\varphi_0}$ est un ss-groupe
de $G$ sss il est formé d'éléments $2:2$ permutables --
auquel cas c'est un sous-groupe commutatif, \deleted{\unsure{tel que}} \inserted{Dans ce}
% archive sheet number: 226

